\documentclass[a4paper]{article}

\usepackage{amsmath}
\usepackage{hyperref}
\usepackage{xcolor}

\newcommand{\entropy}[1]{\mathcal{H}(#1)}
\newcommand{\inst}[1]{\texttt{#1}}

\newcommand{\cmnt}[2]{{\color{red}{\textbf{#1:} #2}}}
\newcommand{\AL}[1]{\cmnt{AL}{#1}}

\newcommand{\code}[1]{\texttt{#1}}

\begin{document}

The following is a discussion of how we handle low-level LLVM operations. The
semantics and grouping are taken from the online reference
manual~\cite{llvm-ref}.

\section*{Uncertainty Coefficient}
We aim to compute the uncertainty coefficient~\cite{press2007numerical} for
each LLVM IR instruction, in order to later use it to provide a program-level
view of the information flow in the program. There is one limitation we will
impose on our analysis, and one assumption. First, we shall limit ourselves to
a stateless view of the program. That is, we provide a static analysis, and
(mostly) do not care what the actual possible inputs to an LLVM instruction
are. If a binary operation takes a pair of registers as inputs, we shall assume
the values are unconstrained, even if at some point in the code, for example,
one of the variables associated with a register was constrained to a constant.
Furthermore, we shall assume all inputs are uniformly distributed. This goes
hand in hand with the previous limitation, and also simplifies the computation.
As such, for some LLVM IR instruction \texttt{inst}, which takes some input
\texttt{I} and returns some output \texttt{O}, we can define the
entropy~(\autoref{eq:squeez}) and the uncertainty
coefficient~(\autoref{eq:unc_coef}) as~\cite{clark2012squeeziness}:

\begin{equation}\label{eq:squeez}
    \entropy{I|O} = \entropy{I} - \entropy{O}
\end{equation}

\begin{equation}\label{eq:unc_coef}
    U(I|O) = \frac{\entropy{I} - \entropy{I|O}}{\entropy{I}} = 1 -
    \frac{\entropy{I|O}}{\entropy{I}} = 1 - \frac{\entropy{I} -
    \entropy{O}}{\entropy{I}} = \frac{\entropy{O}}{\entropy{I}}
\end{equation}

With this as basis, we now turn our attention to the available LLVM IR
instructions, and how we can apply the above formulas to them.

\section*{Computation versus Control Flow}
As discussed above, we are primarily interested in looking at the inputs and
outputs of a particular LLVM IR instruction in order to be able to compute its
uncertainty coefficient. However, first of all, not all instructions have
inputs (e.g. \texttt{ret void}, \texttt{unreachable}), nor do they return some
output (e.g. \texttt{br label <dest>}. Nevertheless, we need to consider how
all instructions might affect information flow in the program.

Thus, we distinguish two overall classes of instructions (outside the taxonomy
in the reference manual itself): computation instructions, those which do
return a value, and thus have an uncertainty coefficient value computable
via~\autoref{eq:unc_coef}, and control flow instructions, which (generally) do
not return a value, and affect the control flow of the program.

\subsection{Computation instructions}\label{sec:comp_inst}
These are instructions that take a number of arguments as inputs and yield one
output. With the assumption that the inputs are uniformly distributed, we can
compute the input entropy as \(\entropy{I} = log_2(i)\), where \(i\) is the
total number of inputs. Thus, for a unary operation, \(i\) would be the total
number of representable bit patterns of size \(b\), where \(b\) is given by the
type encoded in the instruction. With multiple inputs, we duplicate the
representable domain as many times as required, and compute the cartesian
product.

We can thus make our first observation: the maximum uncertainty coefficient of
an operation with \(n\) unconstrained input arguments will be \(1/n\), as the
maximal representable input domain would be \(\mathcal{D}_I \times
\mathcal{D}_I \times ... \times \mathcal{D}_I\), with \(n\) multiplications,
while the output domain is just \(\mathcal{D}_I\). The domain sizes are equal,
as the types of both inputs and output are equal across all LLVM IR
instructions. \AL{Double check}

\subsection{Control flow instructions}
Control flow instructions are not subject to the usual input to output ratio
calculation, by the nature of the fact that they do not have outputs. However,
we believe the information flow in the program is affected by such operations.
Consider a conditional \inst{br} instruction, where the condition is always
false. This makes the code in the target unreachable via that instruction, and
if no other other instruction jumps to it, then that is dead code. By this
measure, we believe we can evaluate the uncertainty coefficient of a control
flow operation.

Let's assume that there are \(n\) states entering the conditional instruction.
Suppose the condition has a \(p\) probability of holding, assuming uniform
input distribution. Thus, \(p*n\) states will follow the \texttt{true} branch,
with \(1-p*n\) in the \texttt{false} branch. By considering these states as
outputs of the instruction, we can then apply our usual entropy calculations.
Thus, the ``output entropy'' over the \texttt{true} branch will be
\(log_2(p*n)\), and over the \texttt{false} branch \(log_2(1-p*n)\). This will
allow us to compute the uncertainty coefficient as normal.

Now, this does mean we need to record some state. The \inst{br} instruction,
for example, in its conditional form, (normally) takes a register as the
conditional argument. We expect that the register is assigned to very closely
to it being used in a conditional instruction, therefore we can record
instructions likely to be used as conditional predicates, such as \inst{icmp}
and \inst{fcmp}, and note the target register. Therefore, the idea is, for each
boolean-returning instruction, we compute the probability of it returning
\texttt{true} or \texttt{false}, by exhaustive evaluation, assuming uniformly
distributed inputs. Then, when we observe a instruction using a recorded
boolean variable, we can use the probability to compute the expected output
state proportions, and assign that as the uncertainty coefficient.

\section{Computing the uncertainty coefficient}
While mathematically, the uncertainty coefficient is computed via the formula
above, we can short-circuit the logic sometimes by considering the semantics of
certain instructions. Thus, we distinguish three further categories of LLVM IR
instructions.

\subsubsection*{Set instructions}
These instructions have a set value for their uncertainty coefficient. For
example, \inst{unreachable} has no defined semantics, aside from saying this
piece of code is unreachable. As such, we set its information flow to 0, as the
expectation is no information should flow beyond this point. Another example is
\inst{fcmp false} (and conversely, \inst{fcmp true}). While both these
instructions take arguments, and the overall semantics intends to perform a
comparison over those arguments, the conditional keyword determines the type of
the comparison. Particularly, \texttt{false} indicates that the operation
always returns \texttt{false}. Therefore, we know that the uncertainty
coefficient is 0, as the output entropy is 0. \AL{Double check}

\subsubsection*{Estimated instructions}
This category of instructions allow us to compute the uncertainty value based
on the arguments passed. For example, \texttt{trunc ... to} takes an argument
of some type, and a secondary type. We can use the bitsizes of the two types in
order to compute the expected entropy loss. Thus, truncating a type of 2 bits
to 1 bit yields an expected uncertainty coefficient value of \(\entropy{O} /
\entropy{I} = log_2{2^1} / log_2{2^2} = 0.5\) \AL{Double check}. This
generalises to \(log_2{2^{bs_o}} / log_2{2^{bs_i}}\).

\subsubsection*{Evaluated instructions}
The final set of instructions are evaluated instructions. For these, we cannot
find a simpler way to compute the uncertainty coefficient than experimental
evaluation. All arithmetic (that is, binary and bitwise binary) operations are
included here, along with a few others (certain comparison predicates).

Since types in LLVM IR are represented by discrete bitsizes, and since large
bit size values yield too big an input space to tractably compute the
uncertainty coefficient values experimentally, we set an arbitrary bit size
limit up to which we perform the evaluation, do an exhaustive evaluation up to
the limit, and then extrapolate the results up to a maximum bit size (in modern
architectures, this would mean 64 bits, albeit 128 bits could be considered).

Regarding the extrapolation, we assume that every LLVM IR operation, supposing
it could have an infinite input space, has some fixed point value as its
corresponding uncertainty coefficient. We further suppose that the experimental
evaluation would approach this fixed point value the higher the bit size, and
that this value cannot be higher than a maximum value (as discussed
in~\autoref{sec:comp_inst}).

We thus perform a naive linear extrapolation process. We compute the difference
between the last two experimentally evaluated bit sizes. We assume, since the
difference between the fixed point value decreases as the bit size increases,
the difference between computed uncertainty coefficients also decreases, thus
making this difference between the last two computed data points the smallest.
Then, we add up this difference as many times as required from the last data
point onwards, in order to reach the maximum bit size under consideration (e.g.
if we evaluated up to a bit size of 20, and would like to extrapolate to a bit
size of 64, we add up 44 times the difference between the computed value for 19
and 20 bits). If this expected maximum bit size value is higher \AL{Lower, for
decaying functions} that the theoretical limit, we simply clamp the maximum
uncertainty coefficient value to the computed limit. Finally, for each discrete
bit size between the last compute bit size and the maximal bit size, we
linearly derive an approximate uncertainty coefficient value.

\section*{Specific instruction discussion}
In the following, we look at specific instructions, split into groups as per
the reference manual.

%\subsection{Unary and Binary operations}
%We combine these together, as the single unary operation available is
%\inst{fneg}. All of these are computation operations. For most of these, we are
%unable to come up with an

\subsection{\inst{add} instruction}\label{sec:discuss_add}
The \inst{add} (and conversely, \inst{sub}) instruction is a good baseline to
compare our uncertainty coefficient analysis against. Assuming no flags are
passed, the basic form of the instruction allows for over- (respectively,
under-) flow. When thinking about how the outputs are distributed, we take note
of two properties. First, as a binary instruction, it takes two operands as
input, and returns one value. Thus, the input space size is \(n^2\), and the
output space size is \(n\), where \(n\) is the total number of representable
values given the type encoded in the instruction. Second, the output is
uniform. That is, each output is present an equal amount of times. \AL{Will
need to discuss a proof here. Probably.}

With these in mind, we conclude that the uncertainty coefficient of \inst{add}
(and \inst{sub}) is the maximum it can be, 0.5.

However, let us now consider the case where one of the operands is constant.
Now, our input space for calculations is reduced to \(n\), the same as the
output space. At the same time, there are no clashes at all in the outputs -
due to wrap-around semantics, the output space is simply returned in its
entirety, with each output value appearing once. Again, we note the uniformity
in the outputs, but since the input space is now smaller, having lost one of
the unbound input variables, the instruction in the form of adding one constant
to a variable does not lose any entropy (i.e. its uncertainty coefficient is
1.0).

Thus, we must consider operations where certain arguments might be set, or have
certain values, as they might greatly affect the computed uncertainty
coefficient in our analysis. At the same time, the presence (or absence) of
certain flags, such as \texttt{nsw} or \texttt{nuw} to yield a \texttt{poison}
value if wrapping occurs also affects the output value. However, in the
interest of time, we currently ignore these flags, as computing the expected
uncertainty values is just extra experimental work, and not much novel
interest.

\subsection{\inst{getelementpointer} instruction}
The form of \inst{getelementpointer} (or \inst{gep} for short) is
\texttt{<result> = getelementptr <ty>, ptr <ptrval>{, <ty> <idx>}*}. That is,
the first type is the type of the given pointer, and indicates how the indexing
should be done. The list of indices are then indices to iteratively dig into
the type, to extract a specific address, corresponding to the object indicated
by the index list.

Normally, looking at the output and considering it is a pointer in stack space,
we could assume the output space is as big as possible stack addresses.
However, the list of indices are encoded in each \inst{gep} call, as well as
the type to interpret the main pointer as. With these two ingredients, if we
consider the pointer offset computed by \inst{gep}, it will be constant, as it
is fully determined by the type and the indices. Thus, a well-defined program
will have no clashes in the output of an instance of a \inst{gep} instruction,
as the output is variable only on the pointer value, and is essentially a fixed
offset from it. This mirrors the constant \inst{add} case
(\autoref{sec:discuss_add}), where we return an output with a constant added to
the input. And similarly, with the input and output spaces equal (the number of
addresses in the stack space), we can conclude that the uncertainty coefficient
is 1.0, i.e. it doesn't lose any entropy.

\subsection{Call instructions}\label{sec:insts_call}
This set of instructions, comprising of \inst{invoke}, \inst{callbr}, and
\inst{call}, all change the control flow to a given function pointer. There are
subtle differences in how to call each of the instructions, but with seemingly
major effects:
\begin{itemize}
    \item [\inst{callbr}] Generally, this should be a rarely appearing
        instruction, used either for a \texttt{goto} or for calling intrinsics.
        The control flow after entering the \inst{callbr} function must then
        continue into one of the provided indirect labels. We shouldn't expect
        this function to appear when compiling normal user-code.
    \item [\inst{call}] This is the usual function calling mechanism we are
        used to: jump and execute the specified instruction, then, on a
        \inst{ret}, come back to the instruction after the \inst{call};
    \item [\inst{invoke}] Similar to \inst{call}, with the main difference that
        a \inst{ret} will not return to the instruction succeeding the
        \inst{invoke}, but rather to the label specified in the arguments of
        the invoke.
\end{itemize}

For the purposes of our analysis, we shall assume that there is no observable
entropy loss at the point of calling, and that any entropy loss calculations
will be done when analysis the code inside the called functions. Of course,
what this means is that if we are calling external functions, then that entropy
loss will not be part of the analysis, but that is a limitation of having a
static analysis tool. A dynamic analyser would be able to ``build up'' entropy
loss as the program executes, by looking at the currently executing
instruction, including from external libraries.

\subsection{\inst{ptrtoaddr} instruction}
This seems to be a fairly new addition to LLVM~\cite{popov2025thisyear}, and is
only slightly different than \inst{ptrtoint}. The difference doesn't seem to
affect functionality \emph{that} much, and it's a very subtle semantic
difference due to the emergence of CHERI, which distinguishes between
provenance and address. At this stage, we handle this instruction similar to
\inst{ptrtoint}, although this seems slightly incorrect\AL{Discuss}.

\subsection{\inst{phi} instruction}
This instruction doesn't perform any computation itself, but rather selects
which value should be returned based on the basic block executed prior to the
current invocation. In terms of inputs / outputs though, this means we have a
list of inputs, with associated basic blocks, and the outputs will be
determined by the prior basic block at runtime. Statically, this means we don't
lose any entropy, as there is a one-to-one mapping for inputs to outputs: the
basic block that was previously executed fully determines the output.

\subsection{Exception handling instruction}
A few instructions deal exclusively with exception handling: \inst{resume},
\inst{catchswitch}, \inst{catchret}, \inst{cleanupret}, \inst{landingpad},
\inst{catchpad}, \inst{cleanuppad}. All of these instructions primarily deal
with handling control flow. Similar to the call
instruction~\ref{sec:insts_call} (and due to time pressure), we shall ignore
analysing the exception handling mechanism, and assume there is no entropy loss
in the functions themselves. There might be information flowing in how the
exceptions and handled which might require special handling in terms of
information flow, but at this stage of the project, there is likely little gain
of handling these particularities compared to analysing the information flow in
the rest of the compiled code.

\section{Dependency Analysis}

Information flow in the program is dependent on how state is consumed across
instructions: some result computed by some instruction will be consumed by
another, and so on until a terminal instruction is reach. Thus, we need a
dependency analysis between LLVM IR instructions to realise this information
flow explicitly, then compute the associated uncertainty coefficient for each
identified path.

Luckily, LLVM already comes with an in-built dependency graph: each
\code{llvm::Instruction} has \code{operands}, which represent the instructions
which are used as input. By starting at terminator instructions (similarly
marked by a boolean \code{isTerminator()} function over
\code{llvm::Instruction}s), we can follow the chain of \code{operands} to get
an explicit dependency graph.

There are two exceptions however: \code{load} and \code{store} instructions. A
\code{load} instruction doesn't have any input operands \AL{double check},
while a \code{store} instruction has both the value to be stored as input, and
the memory location to store to. Thus, we use a separate LLVM analysis tool,
\code{MemorySSA}~\cite{llvm-memoryssa}. With this, additional ``Memory Nodes''
in SSA form are introduced, and we can follow the memory dependency along them
statically, allowing us to model information flow across memory reads and
writes.

With these two approaches, following the syntactical dependency graph, and
following memory reads and associated clobbers via \code{MemorySSA} analysis,
we can model how information flows along the program. Of course, we need to
manually handle some usual issues, such as cycles and path explosion \AL{Which
I haven't gotten to yet ):}.

\bibliographystyle{plain}
\bibliography{notes}

\end{document}
